% Options for packages loaded elsewhere
\PassOptionsToPackage{unicode}{hyperref}
\PassOptionsToPackage{hyphens}{url}
\documentclass[
  a4paper,
]{ltjsarticle}
\usepackage{xcolor}
\usepackage[margin=25mm]{geometry}
\usepackage{amsmath,amssymb}
\setcounter{secnumdepth}{-\maxdimen} % remove section numbering
\usepackage{iftex}
\ifPDFTeX
  \usepackage[T1]{fontenc}
  \usepackage[utf8]{inputenc}
  \usepackage{textcomp} % provide euro and other symbols
\else % if luatex or xetex
  \usepackage{unicode-math} % this also loads fontspec
  \defaultfontfeatures{Scale=MatchLowercase}
  \defaultfontfeatures[\rmfamily]{Ligatures=TeX,Scale=1}
\fi
\usepackage{lmodern}
\ifPDFTeX\else
  % xetex/luatex font selection
\fi
% Use upquote if available, for straight quotes in verbatim environments
\IfFileExists{upquote.sty}{\usepackage{upquote}}{}
\IfFileExists{microtype.sty}{% use microtype if available
  \usepackage[]{microtype}
  \UseMicrotypeSet[protrusion]{basicmath} % disable protrusion for tt fonts
}{}
\makeatletter
\@ifundefined{KOMAClassName}{% if non-KOMA class
  \IfFileExists{parskip.sty}{%
    \usepackage{parskip}
  }{% else
    \setlength{\parindent}{0pt}
    \setlength{\parskip}{6pt plus 2pt minus 1pt}}
}{% if KOMA class
  \KOMAoptions{parskip=half}}
\makeatother
\usepackage{longtable,booktabs,array}
\newcounter{none} % for unnumbered tables
\usepackage{calc} % for calculating minipage widths
% Correct order of tables after \paragraph or \subparagraph
\usepackage{etoolbox}
\makeatletter
\patchcmd\longtable{\par}{\if@noskipsec\mbox{}\fi\par}{}{}
\makeatother
% Allow footnotes in longtable head/foot
\IfFileExists{footnotehyper.sty}{\usepackage{footnotehyper}}{\usepackage{footnote}}
\makesavenoteenv{longtable}
\setlength{\emergencystretch}{3em} % prevent overfull lines
\providecommand{\tightlist}{%
  \setlength{\itemsep}{0pt}\setlength{\parskip}{0pt}}
% Glyph map for lualatex (newunicodechar). Maps symbols that may lack font glyphs.
\usepackage{newunicodechar}
\usepackage{amssymb}
\newunicodechar{∎}{\ensuremath{\blacksquare}}
\newunicodechar{✓}{\ensuremath{\checkmark}}
\newunicodechar{✗}{\ensuremath{\times}}
\newunicodechar{⟹}{\ensuremath{\Longrightarrow}}
\newunicodechar{⟺}{\ensuremath{\Longleftrightarrow}}
\newunicodechar{↪}{\ensuremath{\hookrightarrow}}
\newunicodechar{⌊}{\ensuremath{\lfloor}}
\newunicodechar{⌋}{\ensuremath{\rfloor}}
\newunicodechar{≃}{\ensuremath{\simeq}}
\newunicodechar{≈}{\ensuremath{\approx}}
\newunicodechar{≤}{\ensuremath{\leq}}
\newunicodechar{≥}{\ensuremath{\geq}}
\newunicodechar{≠}{\ensuremath{\neq}}
\newunicodechar{→}{\ensuremath{\rightarrow}}
\newunicodechar{←}{\ensuremath{\leftarrow}}
\newunicodechar{↔}{\ensuremath{\leftrightarrow}}
\newunicodechar{⊥}{\ensuremath{\perp}}
\newunicodechar{√}{\ensuremath{\surd}}
\newunicodechar{∑}{\ensuremath{\sum}}
\newunicodechar{∏}{\ensuremath{\prod}}
\newunicodechar{∞}{\ensuremath{\infty}}
\newunicodechar{∈}{\ensuremath{\in}}
\newunicodechar{∀}{\ensuremath{\forall}}
\newunicodechar{∣}{\ensuremath{\mid}}
\newunicodechar{γ}{\ensuremath{\gamma}}
\newunicodechar{λ}{\ensuremath{\lambda}}
\newunicodechar{μ}{\ensuremath{\mu}}
\newunicodechar{ρ}{\ensuremath{\rho}}
\newunicodechar{σ}{\ensuremath{\sigma}}
\newunicodechar{φ}{\ensuremath{\varphi}}
\newunicodechar{ψ}{\ensuremath{\psi}}
\newunicodechar{θ}{\ensuremath{\theta}}
\newunicodechar{Δ}{\ensuremath{\Delta}}
\newunicodechar{Π}{\ensuremath{\Pi}}
\newunicodechar{Λ}{\ensuremath{\Lambda}}
\newunicodechar{τ}{\ensuremath{\tau}}
\newunicodechar{ε}{\ensuremath{\varepsilon}}
\newunicodechar{ω}{\ensuremath{\omega}}
\newunicodechar{π}{\ensuremath{\pi}}
\newunicodechar{ν}{\ensuremath{\nu}}
\newunicodechar{±}{\ensuremath{\pm}}
\newunicodechar{×}{\ensuremath{\times}}
\newunicodechar{·}{\ensuremath{\cdot}}
\usepackage{bookmark}
\IfFileExists{xurl.sty}{\usepackage{xurl}}{} % add URL line breaks if available
\urlstyle{same}
\hypersetup{
  pdftitle={無名性と情報の最小単位------最小ゼロ閉塞単位は2波、S\_N で割った後に残る最小情報単位は何波か（仮説・検討論文）},
  pdfauthor={Noriaki Kihara (WF System Co., Ltd.), ORCID 0009-0004-6753-4020},
  hidelinks,
  pdfcreator={LaTeX via pandoc}}

\title{無名性と情報の最小単位------最小ゼロ閉塞単位は2波、\(S_N\)
で割った後に残る最小情報単位は何波か（仮説・検討論文）}
\author{Noriaki Kihara (WF System Co., Ltd.), ORCID 0009-0004-6753-4020}
\date{2026-09-12}

\begin{document}
\maketitle

\textbf{シリーズ:}
自己無撞着な関係波閉鎖系の基礎（論文9／全10本・検討論文）\\
\textbf{著者:} 木原 範昭（WF System Co., Ltd.）　\textbf{日付:}
2026-09-12\\
\textbf{Version DOI:} 10.5281/zenodo.22729128\\
\textbf{Concept DOI:} 10.5281/zenodo.22729127\\
\textbf{ORCID:} 0009-0004-6753-4020\\
\textbf{Zenodo:} https://zenodo.org/records/22729128\\

\begin{center}\rule{0.5\linewidth}{0.5pt}\end{center}

\subsection{要旨}\label{ux8981ux65e8}

本論文は\textbf{仮説・検討論文}である（新規走行なし、代数的事実＋既存実データ＋文献照合を、主張分類を明示して整理する）。関係波閉鎖系で「何が情報を担えるか」を無名性の観点から検討する。(1)
辺の番号 \(m\) は頂点置換 \(S_N\) のゲージで、「波 \(m\)
番の型」は物理的な問いではない（数学的事実）。(2)
終状態は厳密等モジュラー（相対ばらつき
\(N=12{:}1.17\times10^{-13}\)、\(N=40{:}3.54\times10^{-14}\)）で、12桁丸めでは異なる
\(|z|\)
が1種類のみ------値でも区別できず、同種粒子の識別不可能性が動力学の帰結として現れる（数値確認）。(3)
閉塞 \(\sum z^2=0\) の非自明解は最小2波（\(z_2=\pm i z_1\)）で、1波は
\(z=0\)
しか許さず情報を持てない（数学的事実＝\textbf{最小閉塞単位は2波}）。ただし
\(\pm\) は pair 交換（\(S_2\)）で反転するので単独では \(S_N\) 不変な bit
でない。primitive block サイズは床依存で論文4 の
\(p=\operatorname{spf}(L)\)（\(L=N/\gcd(N,4)\)、\(8\mid N\)
のみ代数的下限 2
に到達）。ゆえに\textbf{情報の担い手は波個体でなく最小ゼロ閉塞ブロック}（Lam--Leung）だが、「2波＝最小閉塞単位」は「2波＝匿名情報単位」を意味しない。\(\mathcal C_k\simeq Q_{k-2}/S_k\)（射影二次超曲面
\(\dim_{\mathbb C}=k-2\)）を解析すると\textbf{閉塞単独では離散匿名種は生じない}（\(k=2\)
は \(S_2\) で1軌道、\(k\ge3\) は連続モジュライ）；離散な不変量は \(\pm\)
でなくブロック次数
\(p=\operatorname{spf}(L)\)（論文4、閉塞公理＋床の離散対称性が量子化）で、連続モジュライを離散型へ量子化する自己無撞着力学が次の問いである（§5）。(4)
波に \(a,b\)
以外を記録する構造（先行研究の帯・毛レジスタ）は閉塞公理・無名性から導かれない\textbf{足場}であり、そこで読まれた型（ボゾン/フェルミオン）は物理的主張として引用しない（判定の格下げ）。分類：数学的事実／数値確認／仮説を区別。

\begin{center}\rule{0.5\linewidth}{0.5pt}\end{center}

\subsection{1. 課題}\label{ux8ab2ux984c}

この系の波は無名（外部インデックスを持たない）である。無名な対象に何の情報が宿り得るか、粒子種のような区別はそもそも作れるかを、代数と既存実データで検討する。本論文は結論を断定せず、何が確定した事実で何が仮説かを分類して固定する。

\subsection{\texorpdfstring{2. 辺の番号は \(S_N\)
ゲージ（数学的事実）}{2. 辺の番号は S\_N ゲージ（数学的事実）}}\label{ux8fbaux306eux756aux53f7ux306f-s_n-ux30b2ux30fcux30b8ux6570ux5b66ux7684ux4e8bux5b9f}

生成子 \(K_{ef}=A_{ef}\sin(\varphi_f-\varphi_e)\) の \(A\) は \(K_N\)
の辺グラフ隣接で、頂点置換 \(S_N\)
に対し力学は同変である。辺の番号付けは頂点ラベルの選び方の産物で、置換すると別の波に移る。したがって「波
\(m\) 番の型」は物理的な問いではなく、\textbf{読めるのは \(S_N\)
不変量}（総和・多重集合・「型 \(X\)
の波が何本あるか」という個数）だけである。読めないのは「どの波が」（同定・追跡）。

これは論文7
の関係読み出し（絶対法線は状態単独から復元不能、読めるのは面どうしの相対位相）と\textbf{同一原理}である------物理的読出しは絶対ラベル（どの波・どの向き）でなく、変換で不変な関係量に宿る。論文8
まで加えると：(i) \textbf{運動学的無名性}＝最初から波番号は \(S_N\)
ゲージ、(ii) \textbf{動力学的無名化}＝終端では \(|z_e|\)
まで等しくなり値による識別も消える、(iii)
\textbf{残るもの}＝位相関係・面の相対配置・閉塞ブロックの型/多重度という関係的不変量。これがシリーズ全体の関係主義（relationalism）である。

\subsection{3.
終状態は値でも区別できない（数値確認）}\label{ux7d42ux72b6ux614bux306fux5024ux3067ux3082ux533aux5225ux3067ux304dux306aux3044ux6570ux5024ux78baux8a8d}

終状態は厳密等モジュラー：相対ばらつき \(N=12\) で
\(1.17\times10^{-13}\)、\(N=40\) で
\(3.54\times10^{-14}\)。12桁丸めでは異なる \(|z|\)
は1種類のみ（\(66\to1\)、\(780\to1\)）。step0 は \(66/66\)・\(780/780\)
と全て異なるので、無名化は\textbf{動力学が進行させる}。すなわち同種粒子の識別不可能性は初期条件でなく動力学の帰結として出ている（論文5の等振幅・論文8の
\(a_\infty=\sqrt{H/M}\) と同じ事実の別表現、実験7の無名性H定理 \(S\to1\)
とも整合）。

\subsection{4.
情報の最小単位は2波（数学的事実＋仮説）}\label{ux60c5ux5831ux306eux6700ux5c0fux5358ux4f4dux306f2ux6ce2ux6570ux5b66ux7684ux4e8bux5b9fux4eeeux8aac}

閉塞 \(\sum z^2=0\) の非自明解：

{\def\LTcaptype{none} % do not increment counter
\begin{longtable}[]{@{}
  >{\raggedright\arraybackslash}p{(\linewidth - 4\tabcolsep) * \real{0.3333}}
  >{\raggedright\arraybackslash}p{(\linewidth - 4\tabcolsep) * \real{0.3333}}
  >{\raggedright\arraybackslash}p{(\linewidth - 4\tabcolsep) * \real{0.3333}}@{}}
\toprule\noalign{}
\begin{minipage}[b]{\linewidth}\raggedright
波数
\end{minipage} & \begin{minipage}[b]{\linewidth}\raggedright
解
\end{minipage} & \begin{minipage}[b]{\linewidth}\raggedright
自由度
\end{minipage} \\
\midrule\noalign{}
\endhead
\bottomrule\noalign{}
\endlastfoot
1 & \(z^2=0\Rightarrow z=0\) のみ & 無し＝情報を持てない \\
2 & \(z_2=\pm i z_1\)（最小非自明閉塞単位） & \(z_1\)
全体（スケール・位相）。符号 \(\pm\) は pair 交換で反転（下記） \\
3 & 立方根型 \(z_k=z\,\omega^k\) ほか & より広い解多様体 \\
\end{longtable}
}

\textbf{最小閉塞単位は2波（数学的事実）}：非零の閉塞ブロックには最低2波が必要で、1波は
\(z=0\) しか許さず情報を持てない。

\textbf{ただし \(\pm\) は単独では \(S_N\)
不変な2値でない（数学的事実）}：\(z_2=+iz_1\) の対を pair
交換（\(S_2\subset S_N\)）すると \(z'_1=z_2,\ z'_2=z_1\) で
\(z'_2=-iz'_1\)------\textbf{pair の2波を匿名にする \(S_2\)
の作用だけで符号が反転する}。ゆえに \(+i\leftrightarrow-i\) は同一
\(S_N\)
軌道で、右旋/左旋のような二値を物理化するには符号を比較できる追加の関係構造（伝播方向・orientation・第三の関係。論文4
の helicity 議論）が要る。すなわち\textbf{閉塞単位の存在 \(\neq\)
匿名な情報 bit の存在}。

\textbf{primitive block サイズは床依存（論文4）}：論文4 の一般則
\(L=N/\gcd(N,4)\)・\(p=\operatorname{spf}(L)\) で、primitive block は
\(p=2,3,5,7,\ldots\)
を取る（\(N=10\Rightarrow L=5\Rightarrow p=5\)、奇素数
\(N\Rightarrow p=N\)、Lam--Leung）。したがって

\[k_{\min}=2\ (\text{代数的絶対下限})\qquad\text{vs}\qquad k_{\rm primitive}=\operatorname{spf}(L)\ (\text{その床で実現する最小単位}),\]

「最小\textbf{可能}単位」と「実際に許される最小単位」は別で、\(8\mid N\)
だけが下限 2 に到達する。

\textbf{帰結（仮説）と次の問い}：情報の担い手は波個体でなく最小閉塞ブロックである（1波は閉塞できず意味を持てない）。ただし上記より「2波＝最小閉塞単位」は「2波＝匿名な情報単位」を意味しない。本当に問うべきは、閉塞解を対称性で割った後にも離散的に異なる型が残る最小
\(k\) である：

\[\mathcal C_k=\Big\{(z_1,\ldots,z_k):\ \textstyle\sum_{j=1}^k z_j^2=0\Big\}\Big/(\text{全体位相・スケール・}S_k)\]

に離散的に異なる型が残る最小
\(k\)（＝\textbf{最小匿名情報単位}）は何波か。\(k_{\min}=2\)
は最小閉塞単位。最小匿名情報単位は §5
で解析する（\(\mathcal C_k\simeq Q_{k-2}/S_k\)）。粒子種の候補（ブロックサイズ・多重度・組み方）は
\(S_N\) 不変な多重集合として定義でき無名性を破らないが、\(\pm\)
単独は不変 bit でない。

\subsection{\texorpdfstring{5.
閉塞モジュライは連続------離散匿名種には追加の自己無撞着構造が要る（\(\mathcal C_k\simeq Q_{k-2}/S_k\)）}{5. 閉塞モジュライは連続------離散匿名種には追加の自己無撞着構造が要る（\textbackslash mathcal C\_k\textbackslash simeq Q\_\{k-2\}/S\_k）}}\label{ux9589ux585eux30e2ux30b8ux30e5ux30e9ux30a4ux306fux9023ux7d9aux96e2ux6563ux533fux540dux7a2eux306bux306fux8ffdux52a0ux306eux81eaux5df1ux7121ux649eux7740ux69cbux9020ux304cux8981ux308bmathcal-c_ksimeq-q_k-2s_k}

§4 の \(\mathcal C_k\)
は追加計算なしで解析できる。位相・スケールをまとめて
\(z\sim cz\)（\(c\in\mathbb C^\times\)）で割ると、\(\sum_{j=1}^k z_j^2=0\)
は \(\mathbb{CP}^{k-1}\) 上の\textbf{射影二次超曲面}

\[Q_{k-2}=\Big\{[z_1:\cdots:z_k]\in\mathbb{CP}^{k-1}:\ \textstyle\sum_j z_j^2=0\Big\},\qquad \dim_{\mathbb C}Q_{k-2}=k-2,\]

になり、\(\mathcal C_k\simeq Q_{k-2}/S_k\)。

\begin{itemize}
\tightlist
\item
  \textbf{\(k=2\)}：\(z_1^2+z_2^2=0\) の射影解は \([1:i],[1:-i]\)
  の2点（\(\dim_{\mathbb C}Q_0=0\)）だが、\(S_2\) が両者を交換するので
  \(Q_0/S_2\)
  は\textbf{1軌道}------\textbf{2波には閉塞はあるが匿名な2値型はない}（§4
  の \(\pm\) 反転の別証明）。
\item
  \textbf{\(k\ge3\)}：\(\dim_{\mathbb C}Q_{k-2}=k-2\ge1\)。\(S_k\)
  は有限群ゆえ商を取っても一般点近傍の連続次元は消えず、\(\dim_{\mathbb C}(Q_{k-2}/S_k)=k-2\)------解空間は依然\textbf{連続モジュライ}である。
\end{itemize}

したがって\textbf{閉塞条件 \(\sum z^2=0\) 単独では、\(S_k\)
で割った後に離散的に異なる有限個の型（粒子種）は生じない}（\(k=2\) は
\(S_2\) で1種に潰れ、\(k\ge3\)
は連続自由度が残る）。ゆえに問いは「何波なら情報を持てるか」から

\[\boxed{\text{連続な閉塞モジュライを、どの自己無撞着な力学が離散型へ量子化するのか}}\]

へ移る。\textbf{論文4 との接続（量子化の担い手）}：論文4
では高対称床の巡回構造から \(p=\operatorname{spf}(N/\gcd(N,4))\)
という\textbf{離散 block size}
が出た。閉塞方程式単独では連続モジュライなのに、\textbf{高対称床の離散対称性を課すと
block size が量子化される}：

\[\text{閉塞公理}\ +\ \text{床の離散対称性}\ \Longrightarrow\ \text{block size }p=2,3,5,7,\ldots\ \text{の量子化}.\]

現段階で確実に離散な不変量は \(\pm\) ではなく\textbf{ブロック次数 \(p\)
そのもの}であり、その内部にさらに粒子種に相当する離散不変量が生じるかは次の問題である。閉塞の外から論文A
のレジスタ（128個の宣言定数）を足すのでなく、離散型は閉塞＋床＋力学的選択則（orientation・winding・topology
等）が\textbf{内部から}生成せねばならない、という方向を本論文は指す。

\subsection{\texorpdfstring{6. 波に \(a,b\)
以外を記録するのは無名性違反（判定の格下げ）}{6. 波に a,b 以外を記録するのは無名性違反（判定の格下げ）}}\label{ux6ce2ux306b-ab-ux4ee5ux5916ux3092ux8a18ux9332ux3059ux308bux306eux306fux7121ux540dux6027ux9055ux53cdux5224ux5b9aux306eux683cux4e0bux3052}

\begin{itemize}
\tightlist
\item
  本系の波は複素数1個（\(z\in\mathbb C^M\)、\(a+ib\) のみ）＝最小記述。
\item
  先行研究（論文A）は波ごとに \(16\times8=128\) 個の複素数（帯 \(k\)・毛
  \(\eta\) レジスタ）を与える。\(N_n=16,N_\eta=8\)
  は\textbf{宣言定数}で、\(N\) からも閉塞公理 \(\sum z^2=0\)
  からも導かれない。波が内部構造で識別可能になる＝無名性が外から壊されている。
\item
  したがって、そのレジスタ内で読み出された型（ボゾン/フェルミオンの区別・「時計はフェルミオンが点ける」・「電荷＝閉塞欠損」）は\textbf{足場の性質であって閉塞公理から導かれた物理ではない}。本シリーズはこれらを物理的主張として引用しない（判定の格下げ）。型を物理として扱うには、閉塞公理と無名性から帯レジスタ相当の自由度を導出する必要がある（現状未達）。
\end{itemize}

\subsection{7.
主張分類・未決}\label{ux4e3bux5f35ux5206ux985eux672aux6c7a}

\begin{itemize}
\tightlist
\item
  \textbf{数学的事実}：\(S_N\) ゲージ（§2）、\(\sum z^2=0\)
  の最小非自明閉塞単位＝2波 \(z_2=\pm i z_1\)（§4。ただし \(\pm\) は
  pair 交換 \(S_2\) で反転し単独では \(S_N\) 不変 bit
  でない）、primitive block サイズ
  \(=\operatorname{spf}(L)\)（\(L=N/\gcd(N,4)\)、\(8\mid N\)
  のみ代数的下限 2 に到達、§4・論文4）、Lam--Leung 分解。
\item
  \textbf{数値確認}：終状態厳密等モジュラー（§3）。
\item
  \textbf{仮説}：情報の担い手＝最小閉塞ブロック（波個体でない）、粒子種＝\(S_N\)
  不変な多重集合（ブロックサイズ・多重度・組み方）（§4）。
\item
  検証済み（追試B）：終状態のブロック分解の検定を全 \(N=3\)--\(40\)
  で実施した（既存 10000step
  終端状態）。終端状態は厳密等モジュラー（偏差
  \(\sim10^{-13}\)）・\(\Sigma z^2\approx0\) 保存だが \textbf{exact
  \(\pm i\) 対はゼロ}（make\_parent 床も同様。最小 \(\varepsilon\) は大
  \(N\) で \(\sim2\times10^{-6}\)
  まで漸減するのみ）。すなわち最小ブロック（\(\pm i\)
  対）分解は\textbf{巡回Fourier／整数K
  高対称床（\(N=8k\)、論文4）に固有}で、データ系列の終端では立たない。ゆえに「終端ブロック分解＝種の目録」は終端状態の性質としては成立せず、ブロック構造（種に相当する型）は高対称床側（論文4）に宿る------本論文の本旨（情報の担い手はブロック、無名性と両立する型は床の閉塞構造）と整合する。
\item
  \textbf{解析的回答（§5・数学的事実）}：閉塞 \(\sum z^2=0\)
  単独では離散匿名種は生じない------\(\mathcal C_k\simeq Q_{k-2}/S_k\)（射影二次超曲面
  \(\dim_{\mathbb C}Q_{k-2}=k-2\)）で、\(k=2\) は \(S_2\)
  で1軌道に潰れ、\(k\ge3\)
  は連続モジュライ（\(\dim_{\mathbb C}=k-2\ge1\)、有限群 \(S_k\)
  で消えない）。離散なのは \(\pm\) でなく\textbf{ブロック次数
  \(p=\operatorname{spf}(N/\gcd(N,4))\)}（論文4）で、閉塞公理＋床の離散対称性が
  block size を量子化する。
\item
  未決／次の研究問題：連続な閉塞モジュライを離散型へ量子化する自己無撞着な力学は何か（力学的選択則・orientation・winding・topology
  等が内部から生成すべき）。ブロック次数 \(p\)
  の内部にさらに粒子種に相当する離散不変量が生じるか。追試Bは exact
  \(\pm i\) 対（\(p=2\) ブロック）が高対称床 \(8\mid N\) 固有で generic
  終端では消えることを示した（種に相当する型は波個体でなくブロック側に宿り、無名性と両立）。
\end{itemize}

\subsection{関連研究（構造的一致・導出元でない）}\label{ux95a2ux9023ux7814ux7a76ux69cbux9020ux7684ux4e00ux81f4ux5c0eux51faux5143ux3067ux306aux3044}

\begin{itemize}
\tightlist
\item
  消失和（Lam--Leung）：ゼロ閉塞ブロックの素サイズ分解。
\item
  同種粒子の識別不可能性（量子統計）：本系では動力学の帰結として現れる。
\end{itemize}

\subsection{参考文献}\label{ux53c2ux8003ux6587ux732e}

\begin{enumerate}
\def\labelenumi{\arabic{enumi}.}
\tightlist
\item
  木原範昭「自己無撞着な関係波閉鎖系におけるインフレーション的急拡大の機構」Concept
  DOI 10.5281/zenodo.22112008.
\item
  木原範昭「シードが粒子を生む条件と分解能の下限（論文A）」Concept DOI
  10.5281/zenodo.21874481.（§5 で参照する帯・毛レジスタの出典）
\item
  T. Y. Lam and K. H. Leung, ``On vanishing sums of roots of unity'', J.
  Algebra 224 (2000) 91.
\end{enumerate}

\subsection{再現}\label{ux518dux73fe}

本論文は仮説・検討論文で新規走行はない。数値確認（終状態等モジュラー）は既存データ
\texttt{N3\_N40\_long10000\_20260905/}（10000 step
スイープ、\texttt{check\_final\_states\_10000\_v1.py}
等）に依拠。終端ブロック分解の N 依存の検定（追試B）は
\texttt{位相ロック全N確認\_相対平衡\_20260912/終端状態ブロック分解\_N依存\_20260912/}（\texttt{analyze\_terminal\_block\_decomposition\_20260912.py}、\texttt{results/terminal\_block\_summary.csv}、README／REPORT／SHA／run\_all；読出しのみ）。閉塞因数分解は論文4のパッケージと共有。代数（\(\sum z^2=0\)
の最小解・Lam--Leung）は本文で完結。検討の一次記録は発見台帳
\texttt{発見台帳\_干渉保存力学\_20260905起点.md}（検討25・26・27）。

\end{document}
